\chapter{产业链分析与深度调研}
\chapterlead{对应子选题“产业链分析”。年报只讲公司自己，而且讲的是过去；行业、对手和最近发生的事，都要现查。}

\section{为什么先搜集、后分析}

v0 版流程没有独立的资料搜集阶段，行业背景由写作会话边写边补，结果是三份研报全部没有通过“外部信息可追溯”这一项检查：
行业规模、竞争对手的数字写得很具体，但说不清出处。模型记忆里的行业数据可能过时，可能记错，也可能是编造的，而读者无从分辨。

v1 版把资料搜集提到最前面，由一个专门的研究助理会话在任何分析之前完成，之后的框架、预测、估值和文稿都建立在它产出的资料包之上。
写作时用到的每个外部数字都要能对回资料包里的一条登记。改版之后，四份研报全部通过了外部信息可追溯检查，也都没有再命中负向条目。
这两批研报的公司不同，不是严格的对照实验，但方向是清楚的。

\section{七类资料}

研究助理先读年报首页、目录和“管理层讨论与分析”的标题，弄清公司卖什么、分哪几块业务、主要对手是谁，然后按七个类别检索（表~\ref{tab:rs-cat}）。
每个问题至少换两三个检索词，中英文都试；搜索引擎找到候选后，必须打开原文核对数字，只看摘要就登记等于没核对。
来源的优先顺序是：交易所与公司公告，政府、统计局与行业协会，知名研究机构与数据商，主流财经媒体；自媒体、论坛、百科和来源不明的转载不用。

\begin{table}[htbp]
\centering
\caption{资料搜集的七个类别}
\label{tab:rs-cat}
\small
\begin{tabularx}{\linewidth}{@{}l L L@{}}
\toprule
\hd{类别} & \hd{搜什么} & \hd{用在哪里} \\
\midrule
最新业绩 & 年报之后的季报、半年报、业绩预告与快报、月度经营数据 & 预测首年要与最新实际业绩衔接 \\
\rowband
行情 & 最近收盘价与日期、总市值、近三个月走势 & 估值的参照，股价必须有来源和日期 \\
一致预期 & 卖方一致预期、评级分布、目标价区间 & 写预期差：我们和市场哪里不同 \\
\rowband
可比公司 & 3～6 家同业的股价、市值、倍数、收入与利润 & 估值交叉验证 \\
行业 & 行业规模、增速、渗透率、价格指数、上游原材料价格 & 驱动假设的外部依据 \\
\rowband
竞争 & 份额与排名、对手的产能与价格动作、新进入者 & 竞争格局判断要有出处 \\
近期事件 & 近三个月的订单、产能、并购、股权变动、监管、政策 & 年报之后发生的变化 \\
\bottomrule
\end{tabularx}
\end{table}

\section{来源登记与校验}

每一条外部信息都登记进 \file{sources.json}：来源编号、标题、发布方、网址、发布日期、访问日期，以及若干条“事实”。
每条事实包括一句摘自原文、不超过 120 字的引文，和引文中的数值、单位、含义。\code{verify-sources} 检查所登记的数值是否原样出现在引文里，
所以研究助理不能自己换算百分比或单位，原文怎么写就怎么登。排名、“龙头”、政策内容这类定性判断也要登记成事实，只是可以不带数值。

正文引用外部内容时，句末加 \code{[S编号]}。数字护栏把来源登记表里的数值当作合法出处，但要求句子带上对应编号；
“龙头”“领先”“市场份额第一”之类的判断，没有编号同样会被拦下。研报附录的“外部资料来源”由排版工具从登记表自动生成。

表~\ref{tab:rs-stats} 是五个工作区的资料登记规模。每家公司 27～32 个调研问题，44～52 个来源，165～229 条带引文的事实。

\begin{table}[htbp]
\centering
\caption{五个工作区的外部资料登记（\file{research/sources.json}）}
\label{tab:rs-stats}
\small
\begin{tabular}{@{}lrrrrrr@{}}
\toprule
\hd{公司} & \hd{调研问题} & \hd{其中补充调研} & \hd{查到} & \hd{检索词} & \hd{来源} & \hd{登记事实} \\
\midrule
贵州茅台 & 32 & 12 & 30 & 64 & 51 & 229 \\
\rowband
宁德时代 & 27 & 10 & 25 & 54 & 48 & 185 \\
招商轮船 & 29 & 12 & 27 & 56 & 52 & 205 \\
\rowband
中际旭创 & 31 & 16 & 29 & 61 & 51 & 184 \\
宇树科技 & 27 & 15 & 26 & 52 & 44 & 165 \\
\bottomrule
\end{tabular}
\end{table}

查不到的问题标“未找到”，但至少要换过两个检索词。未找到的信息不许用记忆补，相关分析要么改成情景处理，要么删掉；
如果它影响核心结论，写进人工复核清单。

\section{从年报和资料推导分析框架}

资料包交到分析师手里之后，分析师先完成财务抽取，再写分析框架 \file{framework.md}。框架回答三个问题：今年最重要的变化是什么、
背后的驱动是什么、明年会不会延续；然后写出三到五条可以验证的投资主线，每条附带验证点；最后列出驱动清单和外部调研清单。
技能说明里没有任何行业模板，框架完全由模型从这家公司的年报和资料包里推出来。

以中际旭创为例，分析师会话在框架里写下的判断是：2025 年营业收入 382.40 亿元、归母净利润 107.97 亿元（年报第 23 页），利润增速约为收入增速的两倍；
收入增长主要来自量，利润弹性主要来自 800G/1.6T 产品占比提升与硅光方案带来的毛利率抬升；是否延续，要看需求端的北美云厂商资本开支、
供给端的光芯片与物料紧缺、价格端的 1.6T 价格谈判，以及客户集中与地缘政治风险。这四件事随后成为四个深度专题。
图~\ref{fig:rs-chain} 把这条推导画成产业链图。

\begin{figure}[htbp]
\centering
\begin{tikzpicture}[node distance=5mm and 8mm,
  box/.style={draw=#1,fill=#1!8,rounded corners=2pt,align=center,font=\scriptsize,inner sep=4pt,text width=2.55cm,minimum height=1.25cm}]
  \node[box=inkl] (cap) {\textbf{需求端}\\北美云厂商资本开支\\ 1.6T 放量 \textcolor{seal}{[S18–S20]}};
  \node[box=inkl,below=of cap] (chip) {\textbf{供给端}\\光芯片等物料紧缺\\ 预付款锁定 \textcolor{seal}{[S19][S27][S28]}};
  \node[box=seal,right=12mm of cap,yshift=-9mm,text width=3.0cm,minimum height=2.2cm] (co) {\textbf{中际旭创}\\[2pt]光通信收发模块\\ 2025 年销量 2,109 万只\\ 毛利率 42.61\%（p13）};
  \node[box=inkl,right=12mm of co,yshift=9mm] (cust) {\textbf{客户}\\前五大客户占比 75.98\%（p26）\\ 境外收入 90.58\%（p24）};
  \node[box=inkl,below=of cust] (tech) {\textbf{技术路线}\\ 800G → 1.6T、硅光\\ CPO 替代可插拔（p19–20）};
  \node[box=gold,below=8mm of co,text width=6.6cm,minimum height=0.9cm] (risk) {\textbf{价格与政策变量}：1.6T 价格谈判 \textcolor{seal}{[S26][S27]}；美国法案提案 \textcolor{seal}{[S24][S31]}};
  \draw[arr] (cap.east) -- (co.west|-cap.east);
  \draw[arr] (chip.east) -- (co.west|-chip.east);
  \draw[arr] (co.east|-cust.west) -- (cust.west);
  \draw[arr,<->] (co.east|-tech.west) -- (tech.west);
  \draw[darr] (risk.north) -- (co.south);
\end{tikzpicture}
\caption{中际旭创分析框架中的产业链与驱动（据分析师会话 \file{framework.md} 整理；页码指年报，\code{[S编号]} 指来源登记表）}
\label{fig:rs-chain}
\end{figure}

研报中的产业链图由排版工具 reportkit 的流程图组件绘制，节点上的数字同样受数字护栏约束。

\section{深度专题调研}

v1 版研报得分仍然不够（74 分和 67 分），盲评指出的问题集中在分析质量：没有预期差，驱动拆解停留在表面，估值缺少交叉验证。
我们的判断是篇幅和深度不够，研报把大量篇幅花在复述财务数据上。v2 版因此增加了深度调研阶段，并把它写成闸门：
框架里至少列出 3 个深度专题，每个专题写一份不少于 1,500 字、至少引用 3 个独立来源的笔记；补充检索至少 8 个问题；
最终文稿正文不少于 15,000 字，其中深度研究部分不少于 60\%。

五家公司的专题都是分析师会话自己从框架里推出来的，表~\ref{tab:rs-topics} 照录了它们的标题。五家共写了 21 份专题笔记，
合计约 13 万字（按字符计），每家 2.4 万到 3.3 万字。

\begin{table}[htbp]
\centering
\caption{五家公司的深度专题（分析师会话自拟）}
\label{tab:rs-topics}
\small
\begin{tabularx}{\linewidth}{@{}l L@{}}
\toprule
\hd{公司} & \hd{专题} \\
\midrule
贵州茅台 & 价格体系；渠道重构；系列酒出清；行业周期；股东回报与估值底 \\
\rowband
宁德时代 & 单位盈利；储能第二曲线；“去宁化”与份额；全球化 \\
招商轮船 & VLCC 运价与利润弹性；VLCC 供给周期；干散货与 VLOC 长约；竞争格局与政策变量 \\
\rowband
中际旭创 & 需求持续性；价格与毛利率；竞争与技术路线；地缘政治与客户集中 \\
宇树科技 & 人形机器人价格战与毛利率底部；需求结构：科研之后谁接力；份额之争；450 元的股价隐含了什么 \\
\bottomrule
\end{tabularx}
\end{table}

加入深度专题之后，同一评分标准下中际旭创从 74 分升到 90 分，宇树科技从 67 分升到 88 分，提升几乎全部来自分析质量一层（第 9 章）。
逐条看，由“不通过”变为“通过”的条目恰好是核心结论、预期差、驱动拆解、估值交叉验证、情景与敏感性、跟踪指标，正是深度专题要求补上的内容。

\begin{note}[专题不是越多越好]
深度调研阶段的成本不低：分析师会话是五个角色中最贵的一个，1005 批次平均耗时 16.7 分钟、折合 7.19 美元，其中相当一部分花在专题检索上。
我们把专题数量下限定为 3 个，上限不设，实际产出是每家 4～5 个。是否值得为第 5 个专题多花钱，需要消融实验来回答，列在第 13 章的计划里。
\end{note}
